\documentclass[aspectratio=169,11pt]{beamer}
\usetheme{default}
\usecolortheme{default}

\setbeamertemplate{navigation symbols}{}
\setbeamertemplate{footline}[frame number]

\usepackage{amsmath}
\usepackage{booktabs}
\usepackage{graphicx}
\newcommand{\etal}{et al.}

\title{Why this substrate?}
\subtitle{An evidence-driven selection for Project 11}
\author{Srivijayesh Venugopal}
\date{Cranfield IRP 2025-26 --- 11 August 2026}

\begin{document}

\begin{frame}
\titlepage
\end{frame}

% ---------------------------------------------------------------------
\begin{frame}{What Project 11 asks for}

The brief calls for a **fluid-based neural computation** substrate. That imposes four non-negotiable requirements:

\vspace{0.4cm}
\begin{enumerate}
    \item \textbf{Continuity:} computation happens over a spatial field, not at discrete nodes.
    \item \textbf{Parallelism:} independent signal streams can evolve simultaneously.
    \item \textbf{Trainability:} some spatial property of the medium acts as the weights.
    \item \textbf{Native nonlinearity:} the physics itself supplies decision/activation behaviour, not just the readout.
\end{enumerate}

\vspace{0.3cm}
These four requirements are the filter through which every candidate was judged.

\end{frame}

% ---------------------------------------------------------------------
\begin{frame}{The actual search path (not a pre-planned grid search)}

The project moved through candidates as evidence killed them or redirected the focus:

\vspace{0.3cm}
\begin{enumerate}
    \item \textbf{Acoustic waves in channels / porous media} --- implemented, but linear and already well-covered in the literature.
    \item \textbf{Thermal / optothermal modulation} --- own null result: the intended tuning knob did not change the output.
    \item \textbf{Recurrent neural chemical reaction networks} --- trained XOR end-to-end, but no spatial extent.
    \item \textbf{DNA strand-displacement winner-take-all} --- experimentally proven, but spaceless and single-use.
    \item \textbf{Reaction--diffusion excitable media (current)} --- spatial, nonlinear, trainable geometry/light field.
\end{enumerate}

\vspace{0.3cm}
\textbf{Point:} BZ/RD was not the first assumption. It is where the constraints pushed the project.

\end{frame}

% ---------------------------------------------------------------------
\begin{frame}{Selection criteria and their weighting}

\begin{table}
\centering
\footnotesize
\begin{tabular}{p{4.0cm}p{2.0cm}p{7.0cm}}
\toprule
\textbf{Criterion} & \textbf{Weight} & \textbf{Rationale} \\
\midrule
Spatial continuity & Hard requirement & Brief says ``fluid-based'' computation, not discrete reservoirs \\
Parallelism & Hard requirement & Multiple independent streams must coexist \\
Native nonlinearity & High & Otherwise the readout/training carries all the computation \\
Trainable spatial weights & High & Geometry or property fields must be the program \\
Simulation tractability & Medium & Must be verifiable within MSc time \\
Experimental realisability & Medium & Room-temperature, cheap chemistry preferred \\
Speed & Low & Slow is acceptable if the physics is right \\
\bottomrule
\end{tabular}
\end{table}

\vspace{0.2cm}
\textbf{Hard requirements} eliminate any substrate that fails them regardless of other merits.

\end{frame}

% ---------------------------------------------------------------------
\begin{frame}{Wave-based candidates}

\textbf{Acoustic waves (free-field / channel / porous)}
\begin{itemize}
    \item Pros: fast, tractable, FDTD solver implemented.
    \item Cons: passive linear physics; nonlinearity must be imported; field is crowded.
    \item Verdict: \textbf{set aside}. Does not satisfy native nonlinearity; limited novelty headroom.
\end{itemize}

\vspace{0.3cm}
\textbf{Thermal / optothermal media}
\begin{itemize}
    \item Pros: continuous field, easy to simulate.
    \item Cons: own null result showed the heat pattern did not modulate the acoustic signal; diffusive slowness scales as $L^2$.
    \item Verdict: \textbf{set aside}. Fails trainability and speed.
\end{itemize}

\vspace{0.2cm}
Both are continuous, but neither has the required built-in nonlinearity.

\end{frame}

% ---------------------------------------------------------------------
\begin{frame}{Chemistry-based candidates}

\textbf{Well-mixed CRNs / recurrent neural CRNs}
\begin{itemize}
    \item Pros: genuinely trainable; solved XOR end-to-end.
    \item Cons: scalar concentrations only --- no spatial extent, no parallelism across streams.
    \item Verdict: \textbf{set aside}. Fails the hard requirement of spatial continuity.
\end{itemize}

\vspace{0.3cm}
\textbf{DNA strand-displacement winner-take-all}
\begin{itemize}
    \item Pros: experimentally demonstrated pattern recognition.
    \item Cons: spaceless, consumable, weak link to fluid mechanics.
    \item Verdict: \textbf{set aside}. Fails spatial continuity.
\end{itemize}

\vspace{0.2cm}
Trainability is excellent, but these are not \emph{spatial} computers.

\end{frame}

% ---------------------------------------------------------------------
\begin{frame}{Elimination decision tree}

\textbf{Start:} all candidates that can in principle compute.

\vspace{0.2cm}
\begin{enumerate}
    \item \textbf{Is it spatially continuous?}
    \begin{itemize}
        \item No $\rightarrow$ CRN, DNA eliminated.
    \end{itemize}
    \item \textbf{Does it have native nonlinearity?}
    \begin{itemize}
        \item No $\rightarrow$ acoustic, thermal eliminated.
    \end{itemize}
    \item \textbf{What remains?} Reaction--diffusion excitable media.
\end{enumerate}

\vspace{0.3cm}
\textbf{RD checks the remaining boxes:}
\begin{itemize}
    \item Strong native nonlinearity: excitation threshold, refractoriness, annihilation.
    \item Trainable weights: wall geometry, light field $\phi(x,y)$, anisotropic diffusion tensor.
    \item Compact verified model: Oregonator / Barkley.
    \item Experimental path: room-temperature photosensitive BZ layers.
\end{itemize}

\end{frame}

% ---------------------------------------------------------------------
\begin{frame}{Why BZ reaction--diffusion, honestly}

It is not because BZ is the fastest or easiest. It is because it is the only candidate left that satisfies the hard requirements:

\vspace{0.3cm}
\begin{table}
\centering
\footnotesize
\begin{tabular}{lcccc}
\toprule
\textbf{Candidate} & \textbf{Continuous} & \textbf{Parallel} & \textbf{Nonlinear} & \textbf{Trainable space} \\
\midrule
Acoustic & Yes & Yes & No & Geometry \\
Thermal & Yes & Yes & No & Own null result \\
Well-mixed CRN & No & No & Yes & Rates \\
DNA WTA & No & No & Yes & Sequences \\
BZ/RD & Yes & Yes & Yes & Geometry / $\phi$ / tensor \\
\bottomrule
\end{tabular}
\end{table}

\vspace{0.3cm}
The admitted cost is speed: chemical waves travel at mm/min. That is accepted as a substrate limitation, not a flaw in the premise.

\end{frame}

% ---------------------------------------------------------------------
\begin{frame}{What each abandoned path contributed}

\begin{itemize}
    \item \textbf{Acoustics:} established the pulse-in / probe-out experimental grammar and the black-box protocol.
    \item \textbf{Thermal:} showed that a continuous field alone is not enough; the tuning knob must actually change the dynamics.
    \item \textbf{CRN XOR:} proved trainability is possible, but also proved that trainability without space is not what Project 11 asks for.
    \item \textbf{DNA WTA:} confirmed that well-mixed chemistry is experimentally real, but consumable and spaceless.
\end{itemize}

\vspace{0.3cm}
\textbf{Each failure sharpened the criteria.} The pivot to RD was not retreat; it was the elimination of everything that could not meet the brief.

\end{frame}

% ---------------------------------------------------------------------
\begin{frame}{What would change our mind?}

If any of the following were true, the selection would reopen:

\vspace{0.3cm}
\begin{enumerate}
    \item A wave-based substrate with demonstrated \emph{native} nonlinearity at the signal amplitude.
    \item A well-mixed chemistry platform with genuine spatial extent and parallel streams.
    \item An experimental realisation showing the BZ phenomenology cannot be reproduced reliably.
\end{enumerate}

\vspace{0.3cm}
Currently, none of these are the case. The RD/BZ choice stands as the most defensible path against the stated requirements.

\end{frame}

% ---------------------------------------------------------------------
\begin{frame}{Questions}

\begin{enumerate}
    \item Does this evidence-driven elimination read as a rigorous selection process?
    \item Should the criteria weighting be made more explicit in the thesis?
    \item Is the honest pivot history acceptable, or should it be reframed further?
    \item Would a quantitative scoring matrix strengthen the argument?
\end{enumerate}

\end{frame}

\end{document}
